\documentclass[12pt]{article}
\usepackage[T1]{fontenc}
\usepackage[utf8]{inputenc}
\usepackage{lmodern}
\usepackage{textcomp}
\usepackage{enumitem}
\usepackage{geometry}
\geometry{margin=1in}
\begin{document}
\begin{center}
Answer Key - Economics - Graduate Level Assessment 8 \
\small (Answers are ordered exactly as in the question paper.)
\end{center}

\section*{Multiple Choice}
\begin{enumerate}[label=\arabic*.]
\item \textbf{Answer:} B
\\ \textit{Options:} A) It is often used to meet emergencies.; B) Its quantity is controlled by the government rather than by market forces.; C) It is backed by a physical commodity like gold or silver.; D) It encourages the accumulation of debt through interest-free loans.; E) It is used for buying military supplies during a war.
\\ \textit{Note:} The danger inherent in the issuance of fiat money lies in the fact that its quantity is controlled by the government, which can lead to inflation and loss of value if the supply is not managed wisely, unlike currencies regulated by market forces that naturally adjust based on demand and supply.
\item \textbf{Answer:} A
\\ \textit{Options:} A) of the principle of comparative advantage.; B) of the existence of fixed resources.; C) of the scarcity of resources.; D) of the law of decreasing costs.; E) of inefficiencies in the economy.
\\ \textit{Note:} The correct answer is based on the principle of comparative advantage, which indicates that as production of one good increases, resources are not perfectly adaptable to producing other goods, leading to increasing opportunity costs and a concave shape of the production possibilities frontier.
\item \textbf{Answer:} C
\\ \textit{Options:} A) the equilibrium is unstable and each party would like to switch strategies; B) the marginal utility curve intersects the marginal cost curve; C) neither party has an incentive to deviate from his or her strategy; D) the equilibrium is stable and each party would like to switch strategies; E) the equilibrium is where the total revenue equals total cost
\\ \textit{Note:} At a Nash equilibrium, each player's strategy is optimal given the strategy of the other player(s), meaning no player can improve their payoff by unilaterally changing their strategy, thus they have no incentive to deviate.
\item \textbf{Answer:} C
\\ \textit{Options:} A) The quality of housing construction would deteriorate; B) More houses could be bought; C) The housing industry would slump; D) Mutual savings banks would become obsolete; E) Deposits in mutual savings banks would increase
\\ \textit{Note:} The housing industry would slump because increased competition between commercial banks and mutual savings banks could lead to higher interest rates on loans, making it more expensive for consumers to finance home purchases and thereby reducing demand for housing construction.
\item \textbf{Answer:} D
\\ \textit{Options:} A) shift the Phillips curve to the left.; B) shift the investment demand curve to the right.; C) shift the money demand curve to the right.; D) shift the Phillips curve to the right.
\\ \textit{Note:} A negative or contractionary supply shock, such as an increase in oil prices, raises production costs and reduces aggregate supply, leading to higher inflation and lower output, which shifts the Phillips curve to the right, indicating a trade-off between higher inflation and higher unemployment.
\end{enumerate}

\section*{True / False}
\begin{enumerate}[label=\arabic*.]
\setcounter{enumi}{5}
\item \textbf{Answer:} False
\\ \textit{Options:} A) True; B) False
\\ \textit{Note:} The statement is false because, in monopolistic competition, firms have some degree of market power due to product differentiation, allowing them to influence the price of their products rather than being price takers.
\item \textbf{Answer:} True
\\ \textit{Options:} A) True; B) False
\\ \textit{Note:} The statement is true because in a perfectly competitive market, individual firms are price takers due to the presence of many sellers offering identical products, resulting in a horizontal demand curve that indicates consumers will only purchase at the market price.
\item \textbf{Answer:} True
\\ \textit{Options:} A) True; B) False
\\ \textit{Note:} The statement is true because, under the assumption of a stationary time series, the best predictor for the next value of a variable is the mean of its past values, which minimizes the mean squared error of the forecast.
\item \textbf{Answer:} False
\\ \textit{Options:} A) True; B) False
\\ \textit{Note:} Continuously compounded returns represent the exponential growth of an investment over time, whereas changes in asset prices are typically measured as discrete, not continuous, fluctuations, thereby distinguishing the two concepts.
\item \textbf{Answer:} False
\\ \textit{Options:} A) True; B) False
\\ \textit{Note:} A monopolist maximizes profit by setting a price above average total cost, resulting in positive economic profit in the long run, contrary to the conditions of perfect competition where price equals average total cost.
\end{enumerate}

\section*{Long Form Response}
\begin{enumerate}[label=\arabic*.]
\setcounter{enumi}{10}
\item \textbf{Answer:} A capital/output ratio of 2.5 indicates that for every dollar of output, $2.50 is invested in capital. When GNP increases from $985 billion to $1,055 billion, representing a rise of $70 billion, the required investment to support this increase can be calculated as $70 billion multiplied by the capital/output ratio of 2.5, resulting in an investment level of $175 billion. Furthermore, with a marginal propensity to consume of 0.5, any increase in GNP will lead to further changes in consumption and, consequently, a multiplier effect on the economy. This multiplier effect implies that the initial increase in GNP could potentially lead to a total change of \$350 billion when factoring in the consumption response, highlighting the interconnectedness of investment and consumption in economic growth.
\\ \textit{Note:} correct because it accurately applies the capital/output ratio to determine the necessary investment for the increase in GNP, illustrating the relationship between output growth and required capital investment, while also recognizing the role of the marginal propensity to consume in influencing further economic changes.
\item \textbf{Answer:} In the given scenario, the commercial bank has $100,000 in deposits and a required reserve ratio of 20%. This means the bank must maintain $20,000 as reserves, leaving it with $80,000 available for lending. Since the bank currently has $65,000 in loans, it can lend an additional $15,000. When this new loan is made, it will enter the money supply, and through the money multiplier effect, the overall money supply could potentially increase by up to $75,000, as each dollar lent can lead to further deposits and loan creation in the banking system. Thus, the bank's lending capacity not only impacts its balance sheet but also has significant implications for the broader economy.
\\ \textit{Note:} correct because it accurately calculates the available lending capacity of the bank after accounting for required reserves and explains how the subsequent loans can amplify the money supply through the money multiplier effect.
\end{enumerate}

\vfill
\noindent\textit{End of Answer Key}
\end{document}